% Hinweis: Die Ueberschrift (\section*{Abkuerzungsverzeichnis}) und der
% TOC-Eintrag werden in main.tex gesetzt - dort nur dann, wenn tatsaechlich
% Abkuerzungen verwendet wurden (automatisches Ausblenden leerer Verzeichnisse).
% Diese Datei enthaelt daher ausschliesslich das acronym-Environment mit den
% Definitionen; es wird immer geladen, damit \ac{...} im Text funktioniert.
\begin{acronym}[WYSIWYG]\itemsep0pt %der Parameter in Klammern sollte die längste Abkürzung sein. Damit wird der Abstand zwischen Abkürzung und Übersetzung festgelegt
  \acro{OC}{FOM Online Campus}
  \acro{WYSIWYG}{What you see is what you get}
  \acro{Beispiel}{Nicht verwendet, taucht nicht im Abkürzungsverzeichnis auf}
\end{acronym}
